\chapter{Endovascular Aneurysm Repair}

\section*{What Is This Test or Procedure?}
Endovascular aneurysm repair (EVAR) is a minimally invasive procedure that excludes an abdominal aortic aneurysm from the circulation using a stent-graft delivered through the femoral arteries.

\section*{Alternative Names}
EVAR; stent-graft repair; endovascular aortic repair

\section*{Principle}
A fabric-covered metal stent-graft is advanced through the femoral arteries and deployed within the aneurysm under fluoroscopic guidance. The graft relines the aorta, excluding the aneurysm sac from systemic pressure and allowing it to shrink.

\section*{History}
EVAR was introduced by Juan Parodi in 1991. It rapidly became the preferred treatment for anatomically suitable infrarenal abdominal aortic aneurysms because of its lower short-term morbidity than open repair.

\section*{Why This Test or Procedure Is Ordered}
Performed for abdominal aortic aneurysms of sufficient size (typically 5.5 cm in men and 5.0 cm in women) or with rapid growth or symptoms, especially in patients at high risk for open surgery.

\section*{Is This a Primary or Secondary Test or Procedure}
Primary alternative to open repair for anatomically suitable aneurysms.

\section*{How This Test or Procedure Is Performed}
Performed under general or regional anesthesia in an operating room or hybrid suite. Bilateral femoral access is obtained, angiography maps the aneurysm, and the stent-graft is deployed with its main body and limbs sealing against the aortic neck and iliac arteries. Completion angiography confirms exclusion.

\section*{What Preparation Is Needed}
Preoperative CT angiography to assess neck anatomy, renal function evaluation, and discontinuation of anticoagulants. Antibiotics may be given.

\section*{Contraindications and Precautions}
Inadequate proximal neck length or severe iliac tortuosity preclude standard EVAR. Precaution: an unstable patient with rupture may be better served by open repair.

\section*{Risks}
Endoleak (persistent flow into the sac), graft migration or kink, access-site bleeding or pseudoaneurysm, renal failure from contrast, embolization, and limb ischemia.

\section*{What Happens After the Test or Procedure}
Overnight or short hospitalization, early ambulation, and surveillance imaging. Patients typically recover faster than after open repair.

\section*{Understanding Your Results}
Postoperative imaging assesses endoleak and sac size. Type I and III endoleaks require intervention; most type II endoleaks are observed.

\section*{Normal Results}
Exclusion of the aneurysm sac with no endoleak, stable or shrinking sac size, and patent graft limbs.

\section*{Follow-up Tests and Procedures}
CT or ultrasound at 1, 6, and 12 months and then annually for life to monitor the sac and detect late endoleak or device failure.

\section*{References}
\begin{enumerate}
  \item MedlinePlus. Endovascular repair of abdominal aortic aneurysm. U.S. National Library of Medicine; 2025.
  \item Current Medical Diagnosis and Treatment. McGraw-Hill; 2025.
\end{enumerate}

\clearpage
